\chapter{Per iniziare}
\label{ch:getting_started}
\index{Getting started}\index{Primo orario}\index{Setup ambiente}

\epigraph{``Il viaggio di mille miglia comincia con un singolo
passo.''}{\textsc{Lao Tzu}, \emph{Tao Te Ching}, LXIV}

\lettrine[lines=3]{Q}{uesto} capitolo \`e un percorso guidato che
porta un nuovo lettore --- a scelta sviluppatore o operatore di
segreteria didattica --- da una macchina vergine al primo orario
generato e visualizzato a schermo. Le istruzioni sono pensate per
essere copia-incollabili: in caso di errore, il riquadro
\emph{Errori comuni} a fondo capitolo \`e il primo posto da
guardare.

\section{Cosa otteniamo}

A fine capitolo avrai:
\begin{itemize}
\item un'installazione locale funzionante di $\pi$Tantum (backend
  FastAPI + frontend SvelteKit);
\item un profilo di test importato (\texttt{small}: 8 classi,
  19 docenti, 6 indirizzi);
\item un orario settimanale generato dalla pipeline CP-SAT
  ``Phase A + Phase B per giorno'' visibile nella vista
  \emph{Schedule};
\item un primo vincolo personalizzato (in linguaggio DSL) creato
  dall'interfaccia tab Vincoli e ri-applicato a una nuova
  generazione.
\end{itemize}

\sidenote{Tempo stimato per chi parte da zero su Windows 11:
30--40 minuti, di cui 10 di download dei pacchetti, 10 di prima
generazione CP-SAT, 10 di esplorazione UI.}

\section{Prerequisiti}

Servono tre cose installate e in \texttt{\$PATH}:

\begin{itemize}
\item \textbf{Python 3.11 o successivi.} Verifica con
  \texttt{python --version}. Su Windows il binario fornito da
  \emph{Microsoft Store} (\texttt{PythonSoftwareFoundation.Python.3.13})
  funziona. Su macOS conviene installare con
  \texttt{brew install python@3.13}; su Linux il pacchetto della
  distribuzione \`e di solito sufficiente.
\item \textbf{Node.js 20 o superiore.} Verifica con
  \texttt{node --version}. Serve per il dev-server SvelteKit
  della UI. Le istruzioni esatte di installazione sono nel
  capitolo~\ref{ch:installazione} (Installazione).
\item \textbf{Git} per clonare il repository e poter incrementare
  il manuale o il codice senza copiare file a mano.
\end{itemize}

\section{Clonazione e dipendenze}

\sidenote{Tutte le istruzioni assumono una shell aperta nella
directory padre dove si vuole materializzare il repo. Le path
nei comandi usano slash forward (\texttt{/}) per compatibilit\`a
ma su Windows funzionano ugualmente.}

\begin{lstlisting}[style=shell]
git clone https://github.com/gborghi/timetable.git
cd timetable

# (opzionale, raccomandato) ambiente virtuale dedicato
python -m venv .venv
# attivazione: PowerShell
.\.venv\Scripts\Activate.ps1
# attivazione: bash/zsh
source .venv/bin/activate

# dipendenze backend Python
pip install -r webui/backend/requirements.txt

# dipendenze frontend Node
cd webui/frontend && npm install && cd ../..
\end{lstlisting}

L'installazione del backend tira dentro \texttt{ortools},
\texttt{fastapi}, \texttt{sqlalchemy}, \texttt{numpy}, \texttt{scipy}
e una manciata di altri pacchetti. Sulla prima esecuzione
\texttt{pip} pu\`o scaricare $\sim$200 MB di binari: \`e normale.

\begin{erroricomunibox}
\textbf{``ortools wheel non trovata''}. Su Python 3.13
\texttt{ortools} 9.15 \`e disponibile per Windows, macOS e
Linux x86\_64. Se la macchina \`e ARM (Apple Silicon, Linux
ARM64) potrebbe servire una build precompilata pi\`u recente
o l'installazione da sorgente. La via pi\`u breve \`e
\texttt{pip install --upgrade ortools}.

\textbf{``Permission denied'' o virtualenv mal montato}.
Su Windows verifica di aver lanciato la PowerShell come
utente normale (non amministratore) e di non aver impostato
\texttt{ExecutionPolicy} restrittiva: il comando
\texttt{Get-ExecutionPolicy} deve restituire almeno
\texttt{RemoteSigned}.
\end{erroricomunibox}

\section{Avvio dei due processi}

$\pi$Tantum gira come due processi collaborativi:

\begin{itemize}
\item il \textbf{backend} FastAPI che serve l'API REST e la logica
  CP-SAT;
\item il \textbf{frontend} SvelteKit che presenta l'interfaccia
  utente in browser.
\end{itemize}

In due tab di terminale separate (o in due pannelli di tmux):

\begin{lstlisting}[style=shell]
# Tab 1 (backend)
cd webui/backend
uvicorn main:app --reload --port 8000

# Tab 2 (frontend)
cd webui/frontend
npm run dev
\end{lstlisting}

Il backend stampa il messaggio
\texttt{INFO: Uvicorn running on http://0.0.0.0:8000}; il
frontend a sua volta annuncia
\texttt{ready in $\sim$300ms} con l'URL locale.

\sidenote{In alternativa, gli script
\texttt{webui/start.sh} (Linux/macOS),
\texttt{webui/start.bat} (Windows cmd) e
\texttt{webui/start.ps1} (PowerShell) avviano entrambi i
processi in un colpo solo. La sezione \emph{Provalo subito}
qui sotto assume questa modalit\`a.}

\begin{esempiobox}[title=Provalo subito]
\begin{lstlisting}[style=shell]
cd webui
./start.ps1   # Windows PowerShell
# oppure
./start.sh    # Linux / macOS / Git Bash
\end{lstlisting}
Se tutto va, dopo $\sim$10 secondi il browser di default
si apre su \texttt{http://localhost:5173}\,. La schermata
mostra il dashboard $\pi$Tantum con la barra di navigazione
in alto e il pannello \emph{Profili} centrale.
\end{esempiobox}

\section{Importare il profilo \texttt{small}}

I sei profili di test (\texttt{small}, \texttt{medium},
\texttt{big}, \texttt{huge}, \texttt{superhuge}, \texttt{mega})
sono pre-confezionati come pickle in
\texttt{engine/scripts/data/<size>/}. L'import li copia
nel database SQLite in uso, ricreando docenti, classi, materie,
assegnamenti e (se richiesto) le lezioni della soluzione
gi\`a calcolata.

Dal dashboard:

\begin{enumerate}
\item Clicca \emph{Importa profilo}.
\item Seleziona \texttt{small} dal men\`u a tendina.
\item Lascia attivo \emph{Importa lezioni dalla soluzione
  pre-calcolata} se vuoi vedere subito un orario senza
  rilanciare il solver. Disattivalo se vuoi partire da capo.
\item Clicca \emph{Esegui import}.
\end{enumerate}

\begin{avvertenzabox}
L'import \emph{azzera} il database corrente. Eventuali dati
inseriti manualmente vengono persi: salva un dump prima se ti
servono. Il file SQLite \`e
\texttt{webui/backend/timetable.db}.

I pickle di test contengono \emph{solo l'anagrafica}:
docenti, classi, indirizzi, assegnamenti e (eventualmente) le
lezioni gi\`a calcolate. \emph{Non} contengono vincoli
personalizzati. Dopo l'import, le tabelle dei vincoli
(indisponibilit\`a, compresenze, vincoli DSL custom) saranno
\textbf{vuote}. Se vuoi riprodurre uno scenario realistico
con vincoli a strati, devi aggiungerli a mano dalle tab
\emph{Docenti} (per le indisponibilit\`a) e \emph{Vincoli}
(per i vincoli logici e DSL).
\end{avvertenzabox}

\begin{biblioparvabox}
\sffamily
Sotto il cofano, l'endpoint
\texttt{POST /api/dataset/import-profile} chiama
\texttt{optimization.import\_engine\_profile()} che a sua volta
delega a tre helper in \texttt{engine\_io.py}:
\texttt{import\_school\_into\_db()} (anagrafica e classi),
\texttt{import\_profs\_into\_db()} (assegnamenti) e ---
opzionalmente ---
\texttt{import\_solution\_into\_db()} (lezioni da
\texttt{solution\_<profile>\_optimized.pkl}).
\end{biblioparvabox}

\section{Visualizzare l'orario}

Vai sulla tab \emph{Schedule} (\texttt{/schedule}). Se hai
importato la soluzione pre-calcolata vedi una griglia
$36\text{ slot} = 6\text{ giorni} \times 6\text{ ore}$, con
una riga per classe, e in ogni cella la lezione: materia, docente,
eventuale aula. La codifica colore (oro/avorio/indaco) segue il
gusto vintage del frontend.

I tre comandi essenziali:

\begin{itemize}
\item \textbf{Filtri rapidi} (top): selezione classe, docente,
  materia. Quando filtri per docente vedi solo gli slot di quel
  docente; \`e la vista comoda per chi prepara la sostituzione di
  un giorno.
\item \textbf{Drag-and-drop}: trascina una lezione da uno slot
  all'altro. Se lo spostamento crea una violazione hard l'UI
  lampeggia in rosso e annulla; se la soluzione resta valida
  lo slot \`e aggiornato istantaneamente, e il pannello
  \emph{Cost} in alto si aggiorna.
\item \textbf{Esporta XLSX} (in basso a destra): genera un file
  Excel impaginato secondo il modulo classico orario delle scuole
  italiane. Comodo per stampare o condividere via email.
\end{itemize}

\section{Lanciare la pipeline da zero}

Ora rifacciamo il giro generando l'orario senza appoggiarsi alla
soluzione pre-calcolata.

\begin{enumerate}
\item Reimporta \texttt{small} con \emph{Importa lezioni}
  disattivato. Il database si svuota delle lezioni; le tabelle
  Docente, Classe, Materia, Assignment restano popolate.
\item Vai su \emph{Optimize} (\texttt{/optimize}).
\item Scegli la pipeline \emph{Phase A + Phase B (giornaliero)}.
  Tieni i parametri di default: \texttt{time\_a=30s},
  \texttt{time\_b=12s}, workers $= 2$.
\item Clicca \emph{Avvia}. L'evento di progresso compare in basso
  a destra: prima il banner \emph{Phase A in corso\ldots}
  (durante la distribuzione settimanale dei conteggi giornalieri),
  poi sei banner \emph{Phase B [day]} che girano uno dopo l'altro.
\item Al termine il banner verde \emph{Pipeline completata}
  porta automaticamente alla tab Schedule con il nuovo orario
  visualizzato.
\end{enumerate}

\sidenote{La distinzione tra \emph{Phase A} e \emph{Phase B}
\`e tipica di $\pi$Tantum: la prima distribuisce, per ogni
cattedra, il numero di ore tra i sei giorni della settimana
(senza fissare ancora le ore esatte); la seconda, dato il
conteggio per giorno, decide quale slot orario assegnare a
ciascuna lezione. \`E la stessa idea della decomposizione
\emph{schedule = day-distribution + day-assignment} usata
storicamente nei timetabling competition.}

Per profili pi\`u grandi (\texttt{huge}, \texttt{mega}) la
pipeline naturale \`e la \emph{decomposizione spettrale} o la
\emph{branch-and-price} con granularit\`a docente; il
capitolo~\ref{ch:benchmarks} offre la matrice completa dei tempi
e il capitolo~\ref{ch:tecniche_di_ottimizzazione_avanzate}
spiega quando ha senso ciascuna combinazione.

\section{Aggiungere un vincolo dalla UI}

Per chiudere il giro inseriamo un vincolo personalizzato e lo
applichiamo. Apri la tab \emph{Vincoli} (\texttt{/constraints}) e
clicca \emph{+ Nuovo vincolo DSL}.

Esempio (incollabile letterale):

\begin{lstlisting}
forall l in lessons where l.subject == Matematica:
    l.day != 6
\end{lstlisting}

Tradotto in italiano: ``per ogni lezione di matematica, il giorno
non sia il sabato''. Premi \emph{Valida} e poi \emph{Salva}.
Il vincolo entra come \texttt{level=hard} nella tabella
\texttt{general\_constraints}.

Torna su Optimize e rilancia la pipeline. Stavolta la Phase A
si accorge del vincolo via DSL hard e ridistribuisce le ore di
matematica nei primi cinque giorni. Se il vincolo \`e
infattibile (per esempio: gi\`a tutti i sabati erano l'unico
spazio disponibile per matematica) il sistema lo segnala e
suggerisce di abbassarlo a \texttt{soft}.

\begin{esempiobox}[title=Provalo subito]
\begin{lstlisting}[style=shell]
# stesso vincolo, lanciato via API:
curl -X POST http://localhost:8000/api/constraints/general \
  -H 'Content-Type: application/json' \
  -d '{
    "name": "no_math_saturday",
    "expression": "forall l in lessons where l.subject == Matematica: l.day != 6",
    "level": "hard"
  }'
\end{lstlisting}
La risposta JSON contiene \texttt{id} del vincolo creato e
\texttt{parsed\_ast\_json} con l'AST validato. Per cancellarlo:
\texttt{DELETE /api/constraints/general/<id>}.
\end{esempiobox}

\section{Diagnostica e prossimi passi}

A questo punto puoi:

\begin{itemize}
\item esplorare la tab \emph{Diagnostica}
  (\texttt{/diagnostics}) che mostra distribuzioni di carico
  giornaliero per classe e per docente, sesto-ore, buchi,
  e l'andamento di obiettivo della Phase A;
\item leggere il capitolo~\ref{ch:dsl_generico_per_i_vincoli}
  (DSL: tutorial passo-passo) per scrivere vincoli pi\`u
  sofisticati;
\item dare un'occhiata al capitolo~\ref{ch:benchmarks} per
  capire quale tecnica scegliere se la tua scuola somiglia a
  uno dei profili \texttt{huge} o \texttt{mega};
\item se vuoi capire come $\pi$Tantum decide internamente,
  apri il capitolo~\ref{ch:metodo_cpsat} (CP-SAT) e quello
  sulla decomposizione spettrale.
\end{itemize}

\fleuron

\begin{erroricomunibox}
\textbf{``Backend non raggiungibile''}. Il dashboard prova a
contattare \texttt{http://localhost:8000/api/health}; se non
risponde, il banner rosso lo segnala. Verifica che il
processo \texttt{uvicorn} sia ancora in esecuzione e che la
porta 8000 non sia occupata da un altro servizio.

\textbf{``Pipeline blocca su Phase A''}. Su \texttt{small}
Phase A termina in 3--5 secondi; oltre 30 secondi sospetta
che il database sia stato importato male (Assignment senza
ore). Reimporta il profilo e riprova.

\textbf{``DSL: parser error''}. Le keyword sono
\emph{case-sensitive}: \texttt{forall}, \texttt{where},
\texttt{in}, \texttt{exists}, \texttt{count} vanno minuscole.
Gli identificatori (Borghi, Matematica, ecc.) sono valori
letterali, senza virgolette.
\end{erroricomunibox}
